\section{Command line and Python library}
\label{sec:cli}

The application, the \file{raftcorr} command and the Python library share one engine, so the same
images and settings give the same results in all three. The command line suits batches and
computers without a display; the library suits your own scripts. With \file{pip} or a source
installation the command is on the path of that Python environment; the Windows installer puts it
on the path when \ui{Add the ``raftcorr'' command to PATH (new terminals)} is ticked, as it is by
default; on a Mac, use the alias of Section~\ref{sec:installing-macos}.

\subsection{The raftcorr command}

\file{raftcorr} without arguments starts the application. \file{raftcorr -{}-version} prints the
version, \file{raftcorr -h} the commands, and \file{raftcorr <command> -h} the options of one
command.

\begin{center}
\small
\begin{tabularx}{\linewidth}{@{}l X@{}}
  \toprule
  Command & What it does \\
  \midrule
  \file{raftcorr gui [PATH]}      & starts the application, opening an image folder or a
                                    \file{.raftcorr} project if given \\
  \file{raftcorr run FOLDER}      & tracks a sequence and writes a data file and/or a project \\
  \file{raftcorr suggest FOLDER}  & measures the displacement with s8 and suggests a model \\
  \file{raftcorr models}          & lists, downloads or imports the model files \\
  \file{raftcorr selftest}        & runs every installed model on a synthetic image pair \\
  \file{raftcorr info}            & prints the versions of RAFTcorr, Python, PyTorch and the
                                    libraries, and the GPU \\
  \bottomrule
\end{tabularx}
\end{center}

\subsubsection{raftcorr run}

\file{raftcorr run FOLDER [options]} tracks the images of \file{FOLDER} (PNG, JPEG, TIFF, BMP),
sorted in natural order with the first one as the reference image, and writes the results: it is
\ui{Run DIC Analysis} without the window.

\begin{center}
\small
\begin{tabularx}{\linewidth}{@{}l l X@{}}
  \toprule
  Option & Default & Meaning \\
  \midrule
  \file{-{}-roi FILE}    & whole image & region of interest as a mask image, white inside; a mask of
    another size is resized to the images \\
  \file{-{}-model M}     & \file{s4} & \file{s1}, \file{s2}, \file{s4}, \file{s8} or a checkpoint
    path; a missing model file is downloaded \\
  \file{-{}-mode MODE}   & \file{accumulative} & \file{accumulative}, \file{incremental},
    \file{every\_n} or \file{custom} (Section~\ref{sec:workflow}) \\
  \file{-{}-every N}     & 10 & with \file{every\_n}: a new reference every $N$ frames \\
  \file{-{}-keys K,K,…}  & none & with \file{custom}: reference frames as 0-based indices, such as
    \file{0,20,45}; 0 is always one, the last image cannot be one \\
  \file{-{}-iters N}     & 20 & iterations of the model \\
  \file{-{}-device D}    & \file{auto} & \file{auto} (the GPU when CUDA is available), \file{cuda},
    \file{cpu} \\
  \file{-{}-smooth S}    & 0 & Gaussian smoothing of the displacement, $\sigma$ in px; 0 is off \\
  \file{-{}-config FILE} & none & run settings as JSON; replaces \file{-{}-model} to \file{-{}-smooth} \\
  \file{-{}-plain-sort}  & off & sort file names character by character, not naturally \\
  \file{-{}-strain}      & off & also compute the strain \\
  \file{-{}-gauge N}     & 41 & virtual strain gauge size in px, odd \\
  \file{-{}-measure M}   & \file{green\_lagrange} & or \file{infinitesimal}, \file{euler\_almansi} \\
  \file{-{}-edge-trim A} & 0.7 & leave strain empty within $A \times (\text{gauge} - 1)/2$\,px of an
    edge, 0 to 1; 0 keeps every point \\
  \file{-{}-out FILE}    & none & data file: \file{.h5}, \file{.hdf5}, \file{.mat} or \file{.npz} \\
  \file{-{}-project FILE}& none & also save a \file{.raftcorr} project with the results \\
  \file{-{}-overwrite}   & off & replace existing output files \\
  \file{-{}-quiet}       & off & hide the information lines; progress, warnings and errors remain \\
  \bottomrule
\end{tabularx}
\end{center}

The data file holds every variable of Table~\ref{tab:export-variables}, with the residual and, after
\file{-{}-strain}, every strain component. It records no physical units (pixel size 1, unit
\file{px}). The other strain settings keep their defaults: uniform weighting, plane fit, output
spacing 1\,px, no pre-fit smoothing, never across cracks or holes. Regions of interest of later
frames are not available here; use the application or \file{frame\_masks} in Python.

Before tracking, \file{run} checks what would fail at the end: an existing output file without
\file{-{}-overwrite}, a wrong \file{-{}-out} extension, an even gauge, or a trim that would empty the
whole region of interest. Without \file{-{}-out} and \file{-{}-project} it tracks, warns and writes
nothing. Its progress lines, such as \file{[ 25.0\%] Frame 2/5 (reference 1): …}, count frames from
1 like the application.

\begin{lstlisting}
raftcorr run images/ --roi roi.png --model s4 --strain --out results.h5 --project specimen.raftcorr
raftcorr run images/ --model s8 --mode incremental --out results.mat
raftcorr run images/ --mode custom --keys 0,40,80 --out results.npz
raftcorr run images/ --roi roi.png --strain --gauge 31 --edge-trim 0 --out results.h5 --overwrite
raftcorr run images/ --config run.json --out results.h5 --quiet
\end{lstlisting}

A \file{-{}-config} file lists only the keys you set; the others keep their defaults and unknown keys
are ignored. The \file{run\_info.config} entry of an exported parameters file has the same layout,
so you can copy the settings of an earlier run from it.

\begin{lstlisting}
{"model": "s4", "iters": 20, "smoothing_sigma": 0.0,
 "tracking": {"mode": "every_n", "every_n": 20}, "tiling": {"memory_fraction": 0.85}}
\end{lstlisting}

\subsubsection{The other commands}
\label{sec:cli-other}

\begin{description}[leftmargin=0.6cm,style=nextline]
  \item[\file{raftcorr suggest FOLDER [-{}-roi FILE] [-{}-device D]}]
    Measures with s8 the displacement of the region of interest between the reference image and
    up to three later images (a quarter of the way, half way, the last), and prints the finest
    model whose reach is at least 1.3 times the largest displacement (its 99th percentile),
    assuming accumulative tracking. When even s8 falls short, it says to use incremental tracking.
    The frames named in its output are 0-based indices.
  \item[\file{raftcorr models [-{}-download [KEY …]] [-{}-import FILE …]}]
    Lists the four models: name, $s$, reach, the path of the installed file or
    \file{NOT INSTALLED}, a summary, and the folders searched. \file{-{}-download} fetches
    the missing files (all, or the keys given: \file{s1 s2 s4 s8}) into the per-user folder;
    \file{-{}-import} installs \file{.pth} files downloaded by hand from the \file{weights-v1}
    release, recognised by their SHA-256 (Section~\ref{sec:installing-models}).
  \item[\file{raftcorr selftest [-{}-device D] [-{}-gui] [-{}-no-models]}]
    Runs every installed model on a 256\,$\times$\,256\,px synthetic speckle pair with a known
    translation and prints, per model, the measured $u$ and $v$, the time, and \file{ok},
    \file{WRONG}, an error or \file{not installed}. \file{-{}-gui} also opens and closes the main
    window off screen and checks the icons and translations; \file{-{}-no-models} skips the models.
  \item[\file{raftcorr info}]
    Prints the versions of RAFTcorr, Python, the operating system, PyTorch, NumPy, SciPy, PySide6,
    h5py, tifffile, Pillow and Matplotlib, whether CUDA is available, and the GPU with its free
    memory. Add its output to a bug report.
\end{description}

\subsubsection{Exit codes}

\begin{center}
\small
\begin{tabularx}{\linewidth}{@{}c X@{}}
  \toprule
  Code & Meaning \\
  \midrule
  0 & success \\
  1 & \file{selftest}: a model failed, none is installed, or the GUI check failed. Any command: an
      unexpected error, printed as a Python traceback (a folder without images, a reference frame
      out of range, an unknown model key) \\
  2 & a wrong option or value; \file{run}: an output file exists without \file{-{}-overwrite}, a
      wrong \file{-{}-out} extension, or invalid strain settings \\
  3 & \file{models}: a download or an import failed \\
  \bottomrule
\end{tabularx}
\end{center}

\subsection{The Python library}

\begin{lstlisting}[language=Python,showstringspaces=false]
import raftcorr

series = raftcorr.track("images/", roi="roi.png", model="s4")   # DisplacementSeries
strain = raftcorr.compute_strain(series)                        # StrainSeries
u = series.u[4]                        # index 4 = frame 5 on screen; px, +right
exx = strain.components["exx"][4]      # NaN outside the region of interest
\end{lstlisting}

The library follows Section~\ref{sec:conventions}: $+u$ right, $+v$ down, in pixels, in reference
coordinates, and index $k$ is image $k$ counted from 0 (frame $k+1$ in the application).

\begin{lstlisting}[language=Python,showstringspaces=false]
raftcorr.track(images, roi=None, model="s4", config=None, frame_masks=None, callbacks=None)
raftcorr.compute_strain(series, config=None, callbacks=None)
\end{lstlisting}

\begin{center}
\small
\begin{tabularx}{\linewidth}{@{}l X@{}}
  \toprule
  Argument & Meaning \\
  \midrule
  \file{images}      & a folder (natural order), a list of image files, an
    \file{raftcorr.io.ImageSequence}, or a list of 2-D arrays of grey levels 0 to 255 \\
  \file{roi}         & a mask image path (white inside), a boolean array of the image size, or
    \file{None} for the whole image \\
  \file{model}       & \file{s1} to \file{s8} or a checkpoint path, ignored when \file{config} is
    given; a missing model file is downloaded \\
  \file{config}      & \file{track}: a \file{RunConfig}; \file{compute\_strain}: a
    \file{StrainConfig} (default settings of the strain window) \\
  \file{frame\_masks} & \file{\{index: mask\}}, regions of interest of later frames in their own
    coordinates: they exclude what that frame leaves out and are the region of interest when it is
    a reference \\
  \file{callbacks}   & progress, stop and log callbacks (below) \\
  \bottomrule
\end{tabularx}
\end{center}

The settings are frozen dataclasses with the defaults of the application:

\begin{center}
\small
\begin{tabularx}{\linewidth}{@{}l X@{}}
  \toprule
  Class & Fields (default) \\
  \midrule
  \file{RunConfig} & \file{model} (\file{"s4"}), \file{device} (\file{"auto"}), \file{iters} (20),
    \file{corr\_mode} (\file{None} automatic, \file{"tile"} on demand, \file{"volume"} stored),
    \file{smoothing\_sigma} (0.0), \file{tracking}, \file{tiling}; \file{save(path)} and
    \file{load(path)} use the JSON of \file{-{}-config} \\
  \file{TrackingConfig} & \file{mode} (\file{"accumulative"}), \file{every\_n} (10),
    \file{key\_frames} (0-based) \\
  \file{TilingConfig} & \file{margin\_px} (\file{None}: the reach), \file{max\_tile\_pixels}
    (\file{None}: from free GPU memory), \file{memory\_fraction} (0.85), \file{blend\_px} (32) \\
  \file{StrainConfig} & \file{measure} (\file{"green\_lagrange"}), \file{gauge\_px} (41),
    \file{weighting} (\file{"uniform"}, or \file{"gaussian"}), \file{poly\_order} (1, or 2),
    \file{step} (1), \file{gaussian\_sigma} (\file{None}: gauge/4), \file{smoothing\_sigma} (0.0),
    \file{connectivity\_aware} (\file{True}), \file{edge\_trim\_alpha} (0.70); without a control in
    the window: \file{min\_coverage} (0.25), \file{condition\_max} (1000) \\
  \bottomrule
\end{tabularx}
\end{center}

\subsubsection{The result objects}

\begin{center}
\small
\begin{tabularx}{\linewidth}{@{}l X@{}}
  \toprule
  Attribute & Content \\
  \midrule
  \multicolumn{2}{@{}l}{\textbf{DisplacementSeries}} \\
  \file{u}, \file{v}  & $(T, h, w)$ float32, px; NaN outside the region of interest, at untracked
    points and in frames not computed; read-only \\
  \file{mask}, \file{done} & $(h, w)$ region of interest of frame 0; $(T)$ computed frames \\
  \file{grid}         & \file{x0}, \file{y0}, \file{width}, \file{height} and the image size;
    \file{grid.coords()} gives the image coordinates \file{X}, \file{Y} \\
  \file{residual}     & $(T, h, w)$ residual in grey levels, or \file{None} \\
  \file{frame\_names}, \file{n\_frames} & image names; $T$ \\
  \file{completed\_frames} & indices of the computed frames \\
  \file{schedule}     & \file{reference\_of(k)}, the reference of frame $k$ \\
  \file{info}         & \file{model\_name}, \file{model\_sha256}, \file{iters}, \file{device},
    \file{config}, \file{seconds}, \file{notes} \\
  \midrule
  \multicolumn{2}{@{}l}{\textbf{StrainSeries}} \\
  \file{components}   & \file{\{name: (T, h\_s, w\_s)\}} with the names of
    Table~\ref{tab:export-variables}; point $(i, j)$ is at image pixel
    $(x_0 + j\,\text{step},\ y_0 + i\,\text{step})$, with $x_0$, $y_0$ from \file{grid} \\
  \file{step}, \file{done} & output spacing in px; frames with strain \\
  \file{settings}, \file{trimmed} & settings used; edge-trimmed points, or \file{None} \\
  \bottomrule
\end{tabularx}
\end{center}

\subsubsection{Progress and stopping}

A callbacks object has three methods: \file{progress(fraction, message)}, \file{should\_stop()},
polled between frames, and \file{log(message, level)} with the levels \file{debug}, \file{info},
\file{warning} and \file{error}. \file{ConsoleCallbacks} prints like the command line;
\file{FunctionCallbacks} wraps your own functions; the default is silent. When
\file{should\_stop()} returns \file{True}, \file{track} stops after the current frame and returns
what it computed: check \file{series.done}; \file{series.info.notes} records the stop.
\file{compute\_strain} and the exporters raise \file{raftcorr.StopRequested} instead.

\begin{lstlisting}[language=Python,showstringspaces=false]
import raftcorr
from raftcorr import FunctionCallbacks, RunConfig, StrainConfig, TrackingConfig

def show(fraction, message):
    print(f"{100 * fraction:5.1f} %  {message}")

cfg = RunConfig(model="s8", tracking=TrackingConfig(mode="every_n", every_n=20))
series = raftcorr.track("images/", roi="roi.png", config=cfg,
                        callbacks=FunctionCallbacks(progress=show))
strain = raftcorr.compute_strain(series, StrainConfig(gauge_px=31, step=2))
\end{lstlisting}

\subsubsection{Files, probes and projects}

\begin{lstlisting}[language=Python,showstringspaces=false]
from pathlib import Path
from raftcorr.export import PhysicalUnits, export_results, read_results
from raftcorr.analysis import LineGeom, Probe, extract_series
from raftcorr.io.project import Project, save_project

units = PhysicalUnits(scale=0.0125, unit="mm", frame_rate=2.0)   # recorded; data stay in px
export_results("results.h5", series, strain, units=units)        # .h5, .mat or .npz
data = read_results("results.h5")                                # arrays and data["metadata"]

ext = Probe(id=1, kind="line", geometry=LineGeom(100, 200, 400, 200),
            label="E1", color="#ef4444")                          # reference-image pixels
ts = extract_series(series, ext, None, "strain")                  # extensometer strain
ts = extract_series(series, ext, "strain_exx", "mean", strain=strain)

save_project("specimen.raftcorr",
             Project(image_folder=str(Path("images").resolve()),
                     image_files=tuple(series.frame_names), series=series, strain=strain))
\end{lstlisting}

\file{export\_results} writes the data files of the \ui{Data} tab, and \file{read\_results} reads
any of them back with the MATLAB axis order undone. \file{extract\_series} reads a probe as the
\ui{Analysis} tab does and returns \file{frames}, \file{values}, \file{valid\_fraction} and a
\file{status} per frame. It takes the field names of the application (\file{disp\_u},
\file{disp\_v}, \file{disp\_magnitude}, \file{strain\_exx}, \ldots) and, as reduction, a statistic
(\file{value}, \file{mean}, \file{median}, \file{max}, \file{min}, \file{std},
\file{valid\_fraction}) or, for a line, a gauge reading (\file{strain}, \file{true\_strain},
\file{elongation}, \file{cod}, \file{cod\_sliding}, \file{cod\_magnitude}). A project saved this way
opens in the application (Section~\ref{sec:projects}); \file{load\_project} reads one back.
